\section{Euler-Lagrange equation: Proof of Theorem \ref{thm:EulerLagrangeEquation}}

This section is devoted to the proof of Theorem~\ref{thm:EulerLagrangeEquation}
which especially involves the derivation of a formula for the first variation.

By $\Wir$ we will denote the set of \underline injective and \underline regular curves in~$\W$.
The set $\Hir[1,\infty]$ is defined accordingly.
First we will need
the following lemma, to guarantee that $E^{(2)}$ is well defined on a
sufficiently small $\W\capH$ neighborhood of the curve $\g$:


% As in Theorem~\ref{thm:EulerLagrangeEquation} we abbreviate
% \[ U_\eps := [-\tfrac12,-\eps]\cup[\eps,\tfrac12]
%    \qquad\text{for }\eps\in[0,\tfrac12). \]

\begin{lemma}[$\Hir \capH$ is open in $\W \capH$]\label{lem:Upsilon}
 For any $\g\in\Hir \rzd \capH \rzd$ there is some $\tau_0=\tau_0(\g)>0$ with
 \begin{equation}\label{eq:Upsilon}
  \Upsilon := \sett{\g+\h}{\h\in H^{1,\infty}\rzd, \norm\dh_{L^\infty}\le\tau_0}\subset\Hir[1,\infty].
 \end{equation}
 Moreover, there is a constant $c=c(\g)>0$ with
 \begin{equation} \label{eq:fundamentalEstimates}
 \min\{|\tg(u+w) - \tg(u)|, d_{\tg}(u+w,u)\}  \geq c |w|, \quad \quad
 \quad|\dtg(u)| \geq c
 \end{equation}
 for all $\tg \in \Upsilon$ and $(u,w) \in U_0$.
\end{lemma}

\begin{proof}
  We first show that $\g$ is bi-Lipschitz. To this end, choose $\delta\in(0,\frac12)$ with
\begin{equation*}
 \left(\int_{B_r(z)} \int_{B_r(0)} \frac
   {|\dg(u+w)-\dg(u)|^{2}}{|w|^{2}} \d w \d u\right)^{1/2} \leq \tfrac12
\end{equation*}
for all $z \in \mathbb R / \mathbb Z$ and all $r \in[0,\delta]$ which gives
\begin{align*}
  &\frac 1 {2r}\int_{B_r (z)}\bigg| \dg(x) - \frac 1 {2r}  \int_{B_r(z)} \dg(y) \d y\bigg| \d x \\
  &\leq \frac  1 {4r^2} \int_{B_r(z)} \int_{B_r(z)}  |\dg(x) - \dg(y)| \d x \d y \\
  &\leq \left(  \frac 1 {4r^2} \int_{B_r(z)} \int_{B_r(z)}
  |\dg(x) - \dg(y)|^{2} \d x \d y\right)^{1/2} \\
  &\leq \left( \int_{B_r(z)}  \int_{B_r(z)}
  \frac{|\dg(x) - \dg(y)|^{2}}{|x-y|^{2}} \d x \d y\right)^{1/2} \\
  &\leq \tfrac 1 2 .
\end{align*}
Since $\abs{\frac 1 {2r}  \int_{B_r(z)} \dg(y)
  \d y} \leq 1$  we deduce that
\begin{equation*}
  \inf_{\substack{a\in \mathbb R^n \\ |a|\leq 1}} \frac 1 {2r} \int_{B_r(z)} \abs{\dg(y)-a} \d y \leq \tfrac 1 2.
\end{equation*}
For $x,y \in \mathbb R / \mathbb Z$ with $|x-y| \leq 2\delta$ let $r:=
\tfrac12|x-y|$ and $z \in \mathbb R / \mathbb Z$ be the midpoint of
the shorter arc between $x$ and $y$. Then
\begin{align*}
  |\g(x)-\g(y)| &= \sup_{\substack{a\in \mathbb R^n \\ |a|\leq 1}}\;
  \int_{B_r(z)}\langle \dg(t),a\rangle \d t\\
 &= \sup_{\substack{a\in \mathbb R^n \\ |a|\leq 1}} \; \int_{B_r(z)}\langle
 \dg(t) ,\dg(t) + (a-\dg(t))\rangle \d t\\
&\geq \left( 1 - \inf_{\substack{a\in \mathbb R^n \\ |a|\leq 1}}
  \frac 1 {2r}\int_{B_{r}(z) }|\dg(t)-a|\d t \right) |x-y| \\
  &\geq\tfrac 12 |x-y|
\end{align*}
for all $x,y \in \mathbb R / \mathbb Z$ with $|x-y|\leq 2\delta$.
Since $\g$ is embedded and
\[ (x,y) \mapsto \frac
{|\g(y)-\g(x)|}{|y-x|} \]
 defines a continuous positive function on
$I_\delta :=\{(x,y) \in (\mathbb R / \mathbb Z)^2: |x-y| \geq 2\delta\}$, we furthermore have
\begin{equation*}
  |\g(x) - \g(y)| \geq \underbrace{\min_{(\tilde x,\tilde y) \in I_\delta}\frac {|\g(\tilde y)-\g(\tilde x)|}{|\tilde y-\tilde x|}}_{>0} |x-y|.
\end{equation*}
for all $(x,y) \in I_\delta$.
Hence, there is a $c_0=c_0(\g)>0$ with
\begin{equation*}
 |\g(x) - \g(x+w)| \geq c_0 |w|
\end{equation*}
for all $w \in [-1/2, 1/2]$.
Lessening $c_0$ if necessary, we can also achieve by regularity
\[ \abs{\dg}\ge c_0 \qquad\text{on }\R/\Z. \]

Letting $\tau_0:=\tfrac12 c_0$ we obtain
for arbitrary $\tg\in \Upsilon$
\begin{align*}
 \abs{\tg(u+w)-\tg(u)} &\ge \abs{\g(u+w)-\g(u)} - \abs{(\tg-\g)(u+w)-(\tg-\g)(u)} \\
 &\ge c_0\abs w-\norm{\dg-\tdg}_{L^\infty}\abs w \\
 &\ge \tfrac12c_0\abs w
\end{align*}
and
\[ \abs{\tdg}\ge\abs{\dg}-\abs{\tdg-\dg}\ge c_0-\norm{\dg-\tdg}_{L^\infty}\ge\tfrac12c_0. \]
From the latter estimate
we deduce by~\eqref{eq:intrinsic} for $u\in\R/\Z$, $w\in[-\tfrac12,\tfrac12]$
\[ d_{\tg}(u+w,u) \ge d_{\tg}(u\pm w,u) \ge \tfrac12c_0\abs w. \]
We have established~\eqref{eq:fundamentalEstimates} which gives~\eqref{eq:Upsilon}.
\end{proof}

We will use the last lemma to prove the following theorem, from which 
Theorem~\ref{thm:EulerLagrangeEquation} will follow quite easily.

\begin{proposition}\label{prop:formula}
  The energy $\E$ is continuously differentiable on $\br{\Wir\capH}\rzd$.
  The derivative of $\E$ at $\g \in {\Wir\capH}$ in direction $\h \in {\Wir\capH}$ exists and
  is given by
  \begin{equation}\label{eq:formulaD}
  \begin{split}
    &\delta\E(\g;\h)  = \\
    &2\lim_{\varepsilon \searrow 0} \iint\limits_{U_\eps} \Bigg\{
    \left(
      \frac 1 {\abs{\D\g}^2}-\frac 1 {d_\g(u+w,u)^2} \right)
      \left\langle \frac{\dg(u)}{\abs{\dg(u)}^2},\dh(u) \right\rangle 
\\
    &\qquad\qquad {}-
      \left( \frac{\sp{\D\g,\D\h}}{\abs{\D\g}^{4}}
	-\frac{ \left.\tfrac{\d}{\d\tau}\right|_{\tau=0}d_{\g+\tau\h}(u+w,u)}
		{d_\g(u+w,u)^{3}} \right)\Bigg\} 
\\ &\hspace{7.6cm} |\dg(u+w)| |\dg(u)|\d w\d u .
  \end{split}
  \end{equation}
\end{proposition}

As $\g$ is absolutely continuous and regular, the derivative 
$\left.\tfrac{\d}{\d\tau}\right|_{\tau=0}d_{\g+\tau\h}(u+w,u)$ is well-defined for almost all
$(u,w) \in U_0$. From~\eqref{eq:intrinsic} we deduce
\begin{multline}\label{eq:dd} 
 \left.\tfrac{\d}{\d\tau}\right|_{\tau=0} d_{\g+\tau\h} (u+w,u) =
 \begin{cases} |w|\int_{0}^1 
    \left\langle \frac {\dg(u+\sigma w)}{|\dg(u+\sigma w)|} , \dh(u+\sigma w) \right\rangle d\sigma
    &\text{if } \length(\g|_{[u,u+w]}) <\frac 1 2 \length(\g),\\
    -|w|\int_{0}^1
    \left\langle  \frac {\dg(u+\sigma w)}{|\dg(u+\sigma w)|} , \dh(u+\sigma w) \right\rangle d\sigma
    &\text{if } \length(\g|_{[u,u+w]}) > \frac 1 2 \length(\g).
    \end{cases}
\end{multline}

To prove Proposition~\ref{prop:formula}, we will first show that the following approximations
of the energy $\E$, in which we cut off the singular part, are continuously differentiable
and provide a formula for the first variation.
For $\varepsilon\in(0,\tfrac12)$ we set
\begin{equation*}
  \Ee(\g) := \iint\limits_{ U_\eps} 
    \br{\frac1{\abs{\D\g}^2}-\frac1{d_\g(u+w,u)^2}}\abs{\dg(u+w)}\abs{\dg(u)}
    \d w\d u.  
\end{equation*}

\begin{lemma} \label{lem:FrechetDifferentiable}
For $\eps\in(0,\tfrac12)$ the functional $\Ee$ is continuously differentiable on the space of all
injective regular curves in $\W \rzd \cap H^{1,\infty}\rzd$.
The directional derivative at $\g\in \Wir \rzd \cap H^{1,\infty} \rzd$ in direction 
$\h\in \W \rzd \cap H^{1,\infty} \rzd$ is given by
\begin{equation}\label{eq:FrechetDifferentiable}
\begin{split}
    &\delta\Ee(\g;\h)  = \\
    & 2\iint\limits_{U_\eps} \Bigg\{
    \left(
      \frac 1 {\abs{\D\g}^2}-\frac 1 {d_\g(u+w,u)^2} \right)
      \left\langle \frac{\dg(u)}{\abs{\dg(u)}^2},\dh(u) \right\rangle 
\\
    &\qquad \qquad {}-
      \left( \frac{\sp{\D\g,\D\h}}{\abs{\D\g}^{4}}
	-\frac{ \left.\tfrac{\d}{\d\tau}\right|_{\tau=0}d_{\g+\tau\h}(u+w,u)}
		{d_\g(u+w,u)^{3}} \right)\Bigg\} 
\\ &\hspace{7.6cm} |\dg(u+w)| |\dg(u)|\d w\d u .
  \end{split}
\end{equation}
\end{lemma}

\begin{proof}
Applying Lemma~\ref{lem:Upsilon}, we obtain an $\W \capH$-neighborhood $Y \subset \Wir \capH$ of $\g$
such that~\eqref{eq:fundamentalEstimates} uniformly holds on $Y$ for any element in $U_\varepsilon$.
The integrand in~\eqref{eq:FrechetDifferentiable} is almost everywhere
the pointwise derivative of the integrand in $\Ee$. Using~\eqref{eq:intrinsic} and~\eqref{eq:dd}, one sees
furthermore that this pointwise derivative is majorized by some $L^1$-function. So,
Lebesgue's Theorem permits to interchange differentiation and integration which,
by a suitable reparametrization, results in~\eqref{eq:FrechetDifferentiable}.

As for continuity of $\Ee$ and $\delta\Ee$ the only difficulty is to treat the intrinsic distance.
Recalling the continuity of the length functional with respect to absolutely continuous curves
we can directly read off from~\eqref{eq:intrinsic} that
the integrand of $\Ee$ defines a continuous operator
$(H_{\mathrm{ir}}^{3/2}\cap H^{1,\infty})(\mathbb R / \mathbb Z, \mathbb R^n) \to L^1(U_{\varepsilon})$.

Since $\gamma$ is regular, for any $u\in\R/\Z$ there are at most two points $w\in[-\tfrac12,\tfrac12]$ satisfying
$\mathscr L(\gamma|_{[u,u+w]}) = \tfrac 1 2 \mathscr L(\gamma)$ which results in a null set in $U_0$.
Additionally using~\eqref{eq:dd}, we see that the integrand of $\delta\Ee$ gives rise to a continuous mapping
$(H_{\mathrm{ir}}^{3/2}\cap H^{1,\infty})(\mathbb R / \mathbb Z, \mathbb R^n) \times H^{1,\infty} (\mathbb R / \mathbb Z, \mathbb R^n)
\to L^1(U_{\varepsilon})$.
Being linear and bounded in the second component, it can be viewed as a continuous mapping from $\W[1,\infty]$
into the linear bounded operators $\Hir\capH\to L^1(U_\eps)$.

Altogether, the integrand of $\Ee$ is a continuously differentiable functional 
$H_{\mathrm{ir}}^{3/2}\capH\to L^1(U_\eps)$.
The statement now follows from the chain rule and the fact that the integration operator
\[ L^1(U_\varepsilon) \to \mathbb R, \qquad
 g \mapsto \iint_{U_\varepsilon} g(u,w) \d u\d w, \]
is continuously differentiable as it is a bounded linear operator.
\end{proof}

Due to the fact that bounded $L^1$-sequences are not uniformly integrable, the approximations $\Ee$ do not even form
a Cauchy sequence in $C^0(\Hir\capH)$.
In order to prove Proposition~\ref{prop:formula}, we state in Lemma~\ref{lem:CauchySequence} below
that $\Ee$ is \emph{nearly} a Cauchy sequence 
in $C^1(X_{\delta})$ for subsets $X_\delta \subset \Wir\capH$, $\delta\ge0$,
which satisfy the following substitute of the uniform integrability property
\begin{equation}\label{eq:uniformintegrability}
 \limsup_{\varepsilon \rightarrow 0}\sup_{\gamma \in X_\delta} 
  \left(\iint_{\mathbb R / \mathbb Z \times [-\eps, \eps]}
  \frac {|\gamma'(u+w) - \gamma'(u)|^2}{w^2} \d w \d u \right) ^{1/2}
\leq \delta.
\end{equation}
The statement involves the Lipschitz constant
\[ \lip_Y E = \sup_{\substack{f,\tilde f\in Y\\f\ne\tilde f}}\frac{\abs{E(f)-E(\tilde f)}}{\norm{f-\tilde f}} \]
for some real-valued functional $E$ and a subset $Y$ contained in its domain.

\begin{lemma} \label{lem:CauchySequence}
 We have
 \begin{equation*}
  \Ee(\gamma) \rightarrow \E(\gamma)
 \end{equation*}
 for all $\gamma \in \Hir \capH$.

 Furthermore, for any $\gn \in \Wir\capH$ there is an open subset $\Upsilon \subset \Wir\capH$ 
 and a constant $C=C(\gn) < \infty$ such that
\begin{equation}
 \limsup_{\varepsilon_1, \varepsilon_2 \rightarrow 0}\lip_{X_\delta\cap\Upsilon} (\E[\eps_1] - \E[\eps_2])
 \le C \delta\label{eq:CauchySequence2}
\end{equation}
for all subsets $X_\delta \subset \W\capH$ satisfying \eqref{eq:uniformintegrability} with $\delta \in[0,1]$.
\end{lemma}

\begin{proof}
 From Lemma~\ref{lem:Upsilon} we get a $\W \capH$-neighborhood
 $\Upsilon \subset \Wir \capH$ of $\gamma$ such that
 \eqref{eq:fundamentalEstimates} holds for all $\g\in\Upsilon$.
 Making $\Upsilon$ smaller if
 necessary, we may also assume the existence of an $\varepsilon_0 >0$ with
 \begin{equation*}
  d_{\g}(u+w,u) = \mathscr L (\g|_{[u,u+w]})
 \end{equation*}
 for all $\g \in \Upsilon$ and  $w \in [-\varepsilon_0, \varepsilon_0]$.

 \renewcommand{\gt}{\tg}
 \renewcommand{\dgt}{\tg'} 
 In order to bring the integrand in the definition of $\E$ and $\Ee$
 in a more convenient form we
 introduce the function
 \[ g(\zeta,\eta,\vartheta,\iota) := \frac{\zeta^{-2}-\eta^{-2}}{\eta^2-\zeta^2}\vartheta\iota \]
 which is Lipschitz continuous and positive on $[\tilde c,\infty)^4$ for any $\tilde c>0$.
 We define for $u\in\R/\Z$, $w\in[-\eps_0,\eps_0]$
 \begin{align*}
  \G[\g]:(u,w) &\mapsto g\br{\abs{\int_0^1\dg(u+\theta_1w)\d\theta_1},\int_0^1\abs{\dg(u+\theta_2w)}\d\theta_2,\abs{\dg(u+w)},\abs{\dg(u)}}.
 \end{align*}
 We have chosen $\Upsilon$ in such a way that the arguments in $\G[]$ are 
 uniformly bounded away from zero.
 Then we decompose the integrand in the definition of $\E$ for $|w| \leq \varepsilon_0$ into
 \begin{align*}
  &\br{\frac1{\abs{\D\g}^2}-\frac1{d_{\g}(u+w,u)^2}}\abs{\dg(u+w)}\abs{\dg(u)}\\
  &=\frac1{w^2}\br{\frac1{\abs{\int_0^1\dg(u+\theta_1w)\d\theta_1}^2}-\frac1{\br{\int_0^1\abs{\dg(u+\theta_2w)}\d\theta_2}^2}}\abs{\dg(u+w)}\abs{\dg(u)}\\
  &=\G[\g](u,w)\frac{\br{\int_0^1 \abs{\dg(u+\theta_2w)}\d\theta_2}^2 - \abs{\int_0^1\dg(u+\theta_1w)\d\theta_1}^2}{\abs w^{2}}  \\
  &=\G[\g](u,w)\frac{\iint_{[0,1]^2} \br{\abs{\dg(u+\theta_1w)}\abs{\dg(u+\theta_2w)} - \sp{\dg(u+\theta_1w),\dg(u+\theta_2w)}} \d\theta_1\d\theta_2}{\abs w^{2}}.
 \end{align*}
 Using $2\abs a\abs b-2\sp{a,b} = \abs{a-b}^2 - \abs{\abs a-\abs b}^2$ for $a,b\in\R^n$ 
 this can be written as
 \begin{align*}
  &\tfrac12 \G[\g](u,w)\frac{\iint_{[0,1]^2}\abs{\dg(u+\theta_1w)-\dg(u+\theta_2w)}^2\d\theta_1\d\theta_2}{w^2} \\
  &\quad{}-\tfrac12 \G[\g](u,w)\frac{\iint_{[0,1]^2}\br{\abs{\dg(u+\theta_1w)}-\abs{\dg(u+\theta_2w)}}^2\d\theta_1\d\theta_2}{w^2}.
 \end{align*}
 We first use this to get
 \begin{align*}
  \E(\gamma) - \Ee(\gamma) \leq C \int_{\mathbb R / \mathbb Z} \int_{-\varepsilon}^\varepsilon
  \G[\g](u,w)\frac{\iint_{[0,1]^2}\abs{\dg(u+\theta_1w)-\dg(u+\theta_2w)}^2\d\theta_1\d\theta_2}{w^2} 
  dw du  \\
  \leq C  \int_{\mathbb R / \mathbb Z} \int_{-\varepsilon}^\varepsilon
  \frac{\iint_{[0,1]^2}\abs{\dg(u+\theta_1w)-\dg(u+\theta_2w)}^2\d\theta_1\d\theta_2}{w^2} 
  dw du 
  \xrightarrow{\varepsilon \rightarrow 0} 0 
 \end{align*}
 which proves the pointwise convergence stated in the lemma.

 Let now $0<\varepsilon_1 < \varepsilon_2 < \varepsilon_0$ and set
 \[ F:= \E[\varepsilon_1] - \E[\varepsilon_2]. \]
 Using the decomposition of the integrand above, we get 
 \begin{align*}
  F(\g) &=  \tfrac12 \int_{\mathbb R / \mathbb Z}
  \int_{\varepsilon_1 < |w| < \varepsilon_2}\G[\g](u,w)\frac{\iint_{[0,1]^2}\abs{\dg(u+\theta_1w)-\dg(u+\theta_2w)}^2\d\theta_1\d\theta_2}{w^2}\d w\d u \\
  &\quad{}-\tfrac12  \int_{\mathbb R / \mathbb Z}
  \int_{\varepsilon_1 < |w| < \varepsilon_2}\G[\g](u,w)\frac{\iint_{[0,1]^2}\br{\abs{\dg(u+\theta_1w)}-\abs{\dg(u+\theta_2w)}}^2\d\theta_1\d\theta_2}{w^2}\d w\d u \\
  &=: \tfrac12 F_1^{(2)}(\g) - \tfrac12 F_2^{(2)}(\g).
 \end{align*}
 To estimate the difference $F(\gt)-F(\g)$ for $\g,\tg\in Y_{0}$, we first consider
 \begin{align*}
  &\abs{\G[\gt](u,w)-\G(u,w)} \\
  &\le C\abs{\abs{\int_0^1\dtg(u+\theta_1w)\d\theta_1}-\abs{\int_0^1\dg(u+\theta_2w)\d\theta_2}}
  \\
  & \quad{}+C\abs{\int_0^1\br{\abs{\dgt(u+\theta w)}-\abs{\dg(u+\theta w)}}\d\theta} \\
  &\quad{}+C\abs{\abs{\dgt(u+w)}-\abs{\dg(u+w)}} +C\abs{\abs{\dgt(u)}-\abs{\dg(u)}} \\
  &\le C\int_0^1\abs{\dgt(u+\theta w)-\dg(u+\theta w)}\d\theta +C\abs{{\dgt(u+w)}-{\dg(u+w)}} +C\abs{{\dgt(u)}-{\dg(u)}} \\
  &\le C\norm{\dgt-\dg}_{L^\infty}.
 \end{align*}
 We arrive at
 \begin{align*}
  &\abs{F_1^{(2)}(\gt)-F_1^{(2)}(\g)} \\
  &\le \iint\limits_{\R/\Z\times[-\eps_2,\eps_2]}\abs{\G[\gt](u,w)-\G(u,w)}\frac{\iint_{[0,1]^2}\abs{\dgt(u+\theta_1w)-\dgt(u+\theta_2w)}^2\d\theta_1\d\theta_2}{w^2}\d w\d u \\
  &\quad{}+\iint\limits_{\R/\Z\times[-\eps_2,\eps_2]}\abs{\G[\g](u,w)}\\
  & \qquad \qquad  \frac{\iint_{[0,1]^2}\abs{\abs{\dgt(u+\theta_1w)-\dgt(u+\theta_2w)}^2-\abs{\dg(u+\theta_1w)-\dg(u+\theta_2w)}^2}\d\theta_1\d\theta_2}{w^2} \\
  & \hspace{11.2cm}\d w\d u \\
  &\le C\norm{\dgt-\dg}_{L^\infty}\iint\limits_{[0,1]^2}\iint\limits_{\R/\Z\times[-\eps_2,\eps_2]}\frac{\abs{\dgt(u+\theta_1w)-\dgt(u+\theta_2w)}^2}{w^2}\d w\d u\d\theta_1\d\theta_2 \\
  &\quad{}+C\norm\dg_{L^{\infty}}\iint\limits_{[0,1]^2}\iint\limits_{\R/\Z\times[-\eps_2,\eps_2]} \\
  & \qquad  \frac{\abs{(\dgt+\dg)(u+\theta_1w)-(\dgt+\dg)(u+\theta_2w)}\abs{(\dgt-\dg)(u+\theta_1w)-(\dgt-\dg)(u+\theta_2w)}}{w^2} \\
  & \hspace{10.2cm}\d w\d u\d\theta_1\d\theta_2 \\
  &\le C\seminorm{\dgt}_{\W[1/2]_{\scriptstyle 2\eps_2}}^2\norm{\dgt-\dg}_{L^\infty}
    + C \norm\dg_{L^{\infty}}\seminorm{\dgt+\dg}_{\W[1/2]_{\scriptstyle 2\eps_2}}\seminorm{\dgt-\dg}_{\W[1/2]_{\scriptstyle 2\eps_2}}
 \end{align*}
 where we set for $\eps\in(0,\tfrac12)$
 \begin{equation*}
  \seminorm{f}_{\W[1/2]_{\scriptstyle\eps}}
  := \left( \iint_{\R/\Z\times[-\eps,\eps]} \frac{\abs{f(u+w)-f(u)}^2}{w^{2}} \d w\d u\right)^{1/2}.
 \end{equation*}
 For the second term we compute
 \begin{align*}
  &\abs{F_2^{(2)}(\gt)-F_2^{(2)}(\g)} \\
  &\le \! \! \iint\limits_{\R/\Z\times[-\eps_2,\eps_2]} \! \! \! \!\abs{\G[\gt](u,w)-\G(u,w)}\frac{\iint_{[0,1]^2} \!\! \br{\abs{\dgt(u+\theta_1w)}-\abs{\dgt(u+\theta_2w)}}^2\d\theta_1\d\theta_2}{w^2}\d w\d u \\
  &\qquad{}+\iint\limits_{\R/\Z\times[-\eps_2,\eps_2]}\abs{\G[\g](u,w)} \cdot\\
  &\qquad \cdot\frac{\iint_{[0,1]^2}\abs{\br{\abs{\dgt(u+\theta_1w)}-\abs{\dgt(u+\theta_2w)}}^2-\br{\abs{\dg(u+\theta_1w)}-\abs{\dg(u+\theta_2w)}}^2}\d\theta_1\d\theta_2}{w^2}\\
  & \hspace{11.2cm} \d w\d u \\
  &\le C\norm{\dgt-\dg}_{L^\infty}\iint\limits_{[0,1]^2}\iint\limits_{\R/\Z\times[-\eps_2,\eps_2]}\frac{\abs{\dgt(u+\theta_1w)-\dgt(u+\theta_2w)}^2}{w^2}\d w\d u\d\theta_1\d\theta_2 \\
  &\quad{}+C\norm\dg_{L^{\infty}}\iint\limits_{[0,1]^2}\iint\limits_{\R/\Z\times[-\eps_2,\eps_2]}\frac{\abs{\br{\abs{\dgt(u+\theta_1w)}-\abs{\dgt(u+\theta_2w)}}+\br{\abs{\dg(u+\theta_1w)}-\abs{\dg(u+\theta_2w)}}}}{\abs w}\cdot \\
  &\qquad\quad{}\cdot\frac{\abs{\br{\abs{\dgt(u+\theta_1w)}-\abs{\dgt(u+\theta_2w)}}-\br{\abs{\dg(u+\theta_1w)}-\abs{\dg(u+\theta_2w)}}}}{\abs w}\d w\d u\d\theta_1\d\theta_2 \\
  &\le C\norm{\dgt-\dg}_{L^\infty}\seminorm{\dgt}_{\W[1/2]_{\scriptstyle 2\eps_2}}^2
      + C\norm\dg_{L^{\infty}}\seminorm{\abs\dgt+\abs\dg}_{\W[1/2]_{\scriptstyle 2\eps_2}}\seminorm{\abs\dgt-\abs\dg}_{\W[1/2]_{\scriptstyle 2\eps_2}}.
 \end{align*}
 Using the chain and product rule for Sobolev spaces and the formula
 \[
  \abs\dgt-\abs\dg 
  = \frac{\sp{\dgt+\dg,\dgt-\dg}}{\abs\dgt+\abs\dg},
 \]
 we obtain, assuming $\eps_{0}<\tfrac14$, for $C>0$ depending on
 the constant from~\eqref{eq:fundamentalEstimates}
 \begin{align*}
  \seminorm{\abs\dgt-\abs\dg}_{\W[1/2]_{\scriptstyle 2\eps_2}} 
  &\le \seminorm{\abs\dgt-\abs\dg}_{\W[1/2](\R/\Z)} \\
  &\le C \norm{\dgt-\dg}_{\W[1/2]\cap L^{\infty}}
  \norm{\dgt+\dg}_{\W[1/2]\cap L^{\infty}}\br{1+\seminorm\tdg_{\W[1/2]_{\scriptstyle 2\eps_2}}+\seminorm\dg_{\W[1/2]_{\scriptstyle 2\eps_2}}}
  \end{align*}
 where
 \[ \norm\cdot_{\W[1/2]\cap L^{\infty}} := \norm\cdot_{\W[1/2]}+\norm\cdot_{L^{\infty}}. \]
 Hence
 \begin{align*}
    \abs{F_2^{(2)}(\gt)-F_2^{(2)}(\g)}
    &\le C \Bigg(\seminorm{\dgt}_{\W[1/2]_{\scriptstyle 2\eps_2}}^2 +
    \norm\dg_{L^{\infty}}\br{\seminorm\tdg_{\W[1/2]_{\scriptstyle 2\eps_2}}+\seminorm\dg_{\W[1/2]_{\scriptstyle 2\eps_2}}}
    \norm{\dgt+\dg}_{\W[1/2]\cap L^{\infty}}\cdot \\
    &\qquad\qquad\qquad\qquad{}\cdot\br{1+\seminorm\tdg_{\W[1/2]_{\scriptstyle 2\eps_2}}+\seminorm\dg_{\W[1/2]_{\scriptstyle 2\eps_2}}}
     \Bigg)
    \|\dgt-\dg\|_{\W[1/2]\cap L^{\infty}}.
 \end{align*}
 The claim follows from
 \[ \limsup_{\eps_{2}\searrow0}\sup_{\tg\in X_{\delta}\cap Y}\seminorm{\dtg}_{\W[1/2]_{\scriptstyle2\eps_{2}}} \le \delta. \]
\end{proof}

\begin{proof}[Proof of Proposition~\ref{prop:formula}]
 In order to prove that \emph{directional derivatives} exist at $\gn\in\Hir\capH$ for all directions $\h\in\W\capH$
 let $\Upsilon$ be as in the proof of Lemma~\ref{lem:CauchySequence} and
 \begin{equation*}
  X_0 := \sett{ \gn + \tau\h}{\tau \in (-1,1)}.
 \end{equation*}
 First we observe that $X_0$ satisfies~\eqref{eq:uniformintegrability}, thus being an admissible set
 for Lemma~\ref{lem:CauchySequence}. Indeed,
 for $\gt := \gn +\tau\h$, $|\tau| \leq 1$,
 \begin{align}
  &\Bigg( \iint_{\mathbb R / \mathbb Z \times [-\eps, \eps]} 
  \frac {|\g_\tau'(u+w) - \g_\tau'(u)|^2}{|w|^2} \d w\d u \Bigg)^{1/2}  \nonumber \\
  &\leq \left( \iint_{\mathbb R / \mathbb Z \times [-\eps, \eps]}
  \frac {|\dg_{0}(u+w) - \dg_{0}(u)|^2}{|w|^2} \d w \d u \right)^{1/2}
  +  \left( \iint_{\mathbb R / \mathbb Z \times [-\eps, \eps]}
  \frac {|\dh(u+w) - \dh(u)|^2}{|w|^2} \d w \d u \right)^{1/2} 
  \nonumber\\
  &\xrightarrow{\eps\searrow0} 0.\label{eq:X0}
 \end{align}
 From this we deduce,
 for
 \[ f_\eps : \tau\mapsto \Ee(\g_{0}+\tau\h), \]
 that
 \begin{align}
  &\abs{f_{\eps_1}'(\tau)-f_{\eps_2}'(\tau)}
  =\abs{\delta\Ee[\eps_1](\g_0+\tau \h;\h)-\delta\Ee[\eps_2](\g_0+\tau \h;\h)} \nonumber\\
  &\le\limsup_{\theta\to0}\abs{\frac{\Ee[\eps_1](\g_0+(\tau+\theta)\h)-\Ee[\eps_1](\g_0+\tau \h)}{\theta}
   - \frac{\Ee[\eps_2](\g_0+(\tau+\theta)\h)-\Ee[\eps_2](\g_0+\tau \h)}{\theta}} \nonumber\\  
  &\le \lip_{X_0\cap\Upsilon}\br{\Ee[\eps_1]-\Ee[\eps_2]}\norm \h_{\W\cap\W[1	,\infty]}\label{eq:f-f}\\
  &\xrightarrow{\eps_1,\eps_2\searrow0}0\qquad\text{by \eqref{eq:CauchySequence2}.}\nonumber
 \end{align}
 As $\Ee\to\E$ pointwise, this proves that $\seq{f_\eps}$
 is a Cauchy sequence in $C^1((-\tau_0, \tau_0))$ converging to $\E(\g_0 + \tau \h)=\lim_{\eps\searrow0}\Ee(\g_0+\tau \h)$ as 
 $ \varepsilon \rightarrow 0$.
 Hence, especially all directional derivatives of $\E$ exist and
 \begin{equation*}
  \delta \E (\g_0;\h) = \lim_{\varepsilon \searrow 0} \delta \Ee (\g_0;\h)
 \end{equation*}
 for all $\g_0 \in \Wir\cap\W[1,\infty]$, $\h \in \W\cap\W[1,\infty]$.
 
  The next step is to establish \emph{G\^ateaux differentiability}. To this end we merely have to show
 $\delta\E(\g_0,\cdot)\in\br{\W\cap\W[1,\infty]}^*$ for $\g_0\in\Wir\cap\W[1,\infty]$. Linearity carries over from $\Ee$.
 In order to prove boundedness we introduce
 \begin{equation}\label{eq:X_delta}
  X_\delta := \{ \g \in \W : \|\g - \gn\|_{\W} \leq \delta\}\qquad\text{for }\delta\in(0,1]
 \end{equation}
 which also satisfies~\eqref{eq:uniformintegrability} as
 for $\g \in X_\delta$ we have
 \begin{equation}\label{eq:Xdelta}
 \begin{split}
  &\Bigg( \iint_{\mathbb R / \mathbb Z \times [-\eps, \eps]}
  \frac {|\dg(u+w) - \dg(u)|^2}{|w|^2} \d w \d u \Bigg)^{1/2} \\
  &\leq \left( \iint_{\mathbb R / \mathbb Z \times [-\eps, \eps]}
  \frac {|\dgn(u+w) - \dgn(u)|^2}{|w|^2} \d w \d u \right)^{1/2}
 +  \delta \xrightarrow{\eps\searrow0} \delta.
 \end{split}  
 \end{equation}
 Now 
 \begin{align*}
  \delta \E(\gn;\h) 
  &= \delta \Ee(\gn;\h) + \delta \E(\gn;\h) - \delta \Ee(\gn;\h)
  \\
  &= \delta \Ee(\gn;\h) + \lim_{\varepsilon_1 \searrow0} 
   \br{\delta \E[\varepsilon_1] (\gn;\h) - \delta \Ee(\gn;\h)}
 \end{align*}
 and thus, arguing as in~\eqref{eq:f-f} and recalling $\delta\Ee(\g_0;\cdot)\in\br{\W\capH}^*$,
\begin{equation*}
\begin{aligned}
  | \delta \E(\gn;\h) |& \leq | \delta \Ee(\gn;\h) |
  +  \underbrace{\limsup_{\eps_1\searrow0}\lip_{X_\delta\cap\Upsilon}\br{\Ee[\eps_1]-\Ee[\eps]}}_{<\infty} \|\h\|_{\W\cap\W[1,\infty]}
  \leq C \|\h\|_{\W\cap\W[1,\infty]}
\end{aligned}
\end{equation*}
 for all $\gn \in \Wir \capH$ and $\h \in \W \capH$. Hence, $\E$ is G\^ateaux differentiable  and the 
differential is given by
\begin{equation*}
 \left(\E\right)'(\gn) = \delta \E(\gn;\cdot)
\end{equation*}
for all $\gn \in \Wir \capH$, $\h \in \W \capH$.

Finally, to see that the differential is \emph{continuous}, let $\sigma >0$ be given and let us choose
$\delta>0$ and $\varepsilon>0$ so small that
\begin{equation*}
 \lip_{X_\delta\cap\Upsilon}\br{\Ee[\eps_1]-\Ee[\eps_2]} \refeq[\le]{CauchySequence2} C \delta \leq \tfrac13 \sigma
\end{equation*}
for all $\varepsilon_1, \varepsilon_2 < \eps$ where $X_\delta$ is as in~\eqref{eq:X_delta}. Then we
have for $\g \in X_\delta \cap\Upsilon$ and any $\h\in\W\capH$
\begin{align*}
 |\delta E^{(2)}(\g;\h) - \delta E^{(2)}(\gn;\h)|
 &\leq |\delta \E(\g;\h) - \delta \Ee(\g;\h)| + 
  |\delta \Ee(\g;\h) - \delta \Ee(\gn;\h)| 
 \\ 
  &\quad + |\delta \Ee(\gn;\h) - \delta \E(\gn;\h)|
 \\
  &\refeq[\le]{f-f}  |\delta \Ee(\g;\h) - \delta \Ee(\gn;\h)| +  \tfrac23 \sigma \norm\h_{\W\capH}.
\end{align*} 
Since $\Ee$ is $C^1$ we deduce that there is an open neighborhood $V\subset X_\delta$
of $\gn$ such that 
\begin{equation*}
 |\delta \Ee(\g;\h) - \delta \Ee(\gn;\h)| \leq \tfrac13\sigma \norm\h_{\W\capH}
\end{equation*}
and hence
\begin{equation*}
  |\delta \E(\g;\h) - \delta \E(\gn;\h)| \leq \sigma \norm\h_{\W\capH}
\end{equation*}
for all $\g\in V$.
This proves that $\left(\E \right)'$ is continuous from $\Wir\capH$ into $\left( \W \capH \right)^\ast$ and
hence $\E$ is $C^1(\Wir$ $\capH)$.
\end{proof}


\begin{proof}[Proof of Theorem~\ref{thm:EulerLagrangeEquation}]
 Using Proposition~\ref{prop:formula} we merely have to derive the formula of the first variation
 for a curve $\g \in \Wir\capH$ parametrized by arc-length
 and $\h \in \W\capH$.
 As $\abs\dg\equiv1$ a.~e.\@ we deduce from Lemma~\ref{lem:FrechetDifferentiable} 
 and Equation~\eqref{eq:dd}
 \begin{align*}
  \delta\E(\g;\h) &\xleftarrow{\varepsilon \searrow 0} 2\iint_{U_\varepsilon}
\Bigg\{\left(\frac 1 {|\g(u+w) - \g(u)|^2} - \frac 1 {w^2} \right)
  \left\langle \dg(u), \dh(u) \right\rangle  \\
  &\qquad {}- 2  \left(\frac{\left\langle \D\g ,\D \h \right\rangle }{|\D\g|^{4}}
    -\frac{\left.\tfrac{\d}{\d\tau}\right|_{\tau=0}d_{\g+\tau\h}(u+w,u)}{|w|^{3}}\right)\Bigg\}
  \d w\d u 
 \end{align*}
 where now 
 \begin{equation*}
  \left.\tfrac{\d}{\d\tau}\right|_{\tau=0}d_{\g+\tau\h}(u+w,u) = |w|\int_0^1 \left\langle \dg(u+ \theta w), \dh(u+\theta w) \right\rangle d\theta
 \end{equation*}
 for all $(u,w) \in \mathbb R / \mathbb Z \times (-\tfrac12, \tfrac12)$. Hence,
 \begin{align*}
  &\delta\E(\g;\h) 
  \xleftarrow{\varepsilon \searrow 0}2\iint_{U_\varepsilon}  \Bigg\{
    \left(\frac 1 {|\g(u+w) - \g(u)|^2} - \frac 1 {w^2} \right) 
    \left\langle \dg(u), \dh(u) \right\rangle \\
    &\quad{} -   \left(\frac{\left\langle \D\g ,\D \h \right\rangle }{|\D\g|^{4}}
    -\frac{\int_0^1 \left\langle \dg(u+\theta w ), \dh(u+\theta w) \right\rangle d\theta}
    {w^2}\right) \Bigg\}\d w\d u \\
  & = 2\iint_{U_\varepsilon} \Bigg\{
    \left(\frac 1 {|\g(u+w) - \g(u)|^2} - \frac 1 {w^2} \right) 
    \left\langle \dg(u), \dh(u) \right\rangle \\
    &\qquad\qquad{} - \left(\frac{\left\langle \D\g ,\D \h \right\rangle }{|\D\g|^{4}}
    -\frac{\left\langle \dg(u), \dh(u) \right\rangle}{w^2}\right) \Bigg\}\d w\d u\\
& = 2\iint_{U_\varepsilon} 
  \br{\frac  {\left\langle \dg(u), \dh(u) \right\rangle} {|\g(u+w) - \g(u)|^2}
            - \frac{\left\langle \D\g ,\D \h \right\rangle }{|\D\g|^{4}} } \d w\d u.
\end{align*}
\end{proof}